\documentclass[12pt,a4paper]{article}

\usepackage[spanish,es-nodecimaldot]{babel}
\usepackage[utf8]{inputenc}
\usepackage[T1]{fontenc}
\usepackage{lmodern}
\usepackage{geometry}
\usepackage{microtype}
\usepackage{setspace}
\usepackage{parskip}
\usepackage{enumitem}
\usepackage{xcolor}
\usepackage{listings}
\usepackage{hyperref}

\geometry{margin=2.5cm}
\onehalfspacing
\setlist[itemize]{leftmargin=*,itemsep=0.2em}
\setlist[enumerate]{leftmargin=*,itemsep=0.2em}

\hypersetup{
  colorlinks=true,
  linkcolor=black,
  urlcolor=blue,
  pdftitle={Resolución del Práctico 2 - Testing de Software 2024},
  pdfauthor={}
}

\lstdefinestyle{codestyle}{
  basicstyle=\ttfamily\small,
  frame=single,
  breaklines=true,
  showstringspaces=false
}

\title{\textbf{Resolución del Práctico 2}\\Testing de Software 2024}
\author{}
\date{}

\begin{document}

\maketitle
\tableofcontents
\newpage

\section*{Introducción}
\addcontentsline{toc}{section}{Introducción}
Este documento presenta la resolución integral del \textit{Práctico 2}. Se mantiene el contenido técnico original y se mejora su presentación con una estructura académica uniforme, redacción formal y organización por ejercicio.

\section{Ejercicio 1}
Se realizó la lectura del capítulo 3 del libro \textit{Introduction to Software Testing} (Ammann y Offutt, segunda edición), conforme a la consigna.

Material de referencia utilizado en el repositorio:
\begin{itemize}
  \item \texttt{material/Introduction to Software Testing.pdf}
\end{itemize}

\section{Ejercicio 2}
Se migraron a tests parametrizados los casos desarrollados en el ejercicio 3 del Práctico 1, ejecutándolos sobre la implementación provista por cátedra:
\begin{itemize}
  \item \texttt{src/main/java/assignment2\_exercises/SimpleRoutines.java}
\end{itemize}

\subsection{Métodos cubiertos}
\begin{itemize}
  \item \texttt{findLast(int[] x, int y)}
  \item \texttt{lastZero(int[] x)}
  \item \texttt{countPositive(int[] x)}
  \item \texttt{oddOrPos(int[] x)}
\end{itemize}

\subsection{Archivos agregados}
\begin{itemize}
  \item \texttt{src/test/java/assignment2\_exercises/SimpleRoutinesParameterizedTest.java}
  \item \texttt{src/test/resources/assignment2\_exercises/odd\_or\_pos\_cases.csv}
\end{itemize}

\subsection{Estrategia de parametrización}
\begin{itemize}
  \item \texttt{@MethodSource} en \texttt{lastZero} y \texttt{countPositive}.
  \item \texttt{@CsvSource} en \texttt{findLast}.
  \item \texttt{@CsvFileSource} en \texttt{oddOrPos}.
\end{itemize}

Con esta combinación se satisface lo solicitado: utilizar métodos proveedores y parámetros en formato CSV.

\subsection{Ejecución}
\begin{lstlisting}[style=codestyle]
mvn -Dmaven.repo.local=.m2 -Djacoco.skip=true -Dtest=SimpleRoutinesParameterizedTest test
\end{lstlisting}

Se utilizó \texttt{-Djacoco.skip=true} debido a incompatibilidades de la versión JaCoCo del template (0.8.2) con JDK actuales.

\subsection{Resultado esperado}
La suite ejecuta 10 casos. Sobre la versión defectuosa de \texttt{SimpleRoutines} fallan los 4 casos que exponen los defectos conocidos:
\begin{itemize}
  \item \texttt{findLast}: no evalúa el índice \texttt{0}.
  \item \texttt{lastZero}: retorna el primer cero en lugar del último.
  \item \texttt{countPositive}: cuenta \texttt{0} como positivo.
  \item \texttt{oddOrPos}: no contabiliza impares negativos.
\end{itemize}

\section{Ejercicio 3}
Se desarrolló una suite de tests para todos los métodos de \texttt{StackAr} con JUnit 5, utilizando fixture compartido con \texttt{@BeforeEach} e identificando explícitamente las fases \textit{arrange}, \textit{act} y \textit{assert}.

\subsection{Archivos trabajados}
\begin{itemize}
  \item \texttt{src/main/java/assignment2\_exercises/stack/StackAr.java}
  \item \texttt{src/test/java/assignment2\_exercises/stack/StackArTest.java}
\end{itemize}

\subsection{Cobertura de pruebas}
Se incluyeron pruebas para:
\begin{itemize}
  \item constructor con capacidad inválida,
  \item \texttt{size},
  \item \texttt{isEmpty},
  \item \texttt{isFull},
  \item \texttt{push},
  \item \texttt{pop},
  \item \texttt{top},
  \item \texttt{makeEmpty},
  \item \texttt{equals},
  \item \texttt{hashCode},
  \item \texttt{toString},
  \item \texttt{repOk}.
\end{itemize}

Además, se agregaron tests negativos adicionales, por ejemplo:
\begin{itemize}
  \item \texttt{push} en pila llena,
  \item \texttt{push(null)},
  \item \texttt{pop} en pila vacía,
  \item \texttt{top} en pila vacía.
\end{itemize}

\subsection{Cambios en la implementación}
\begin{enumerate}
  \item Se completó \texttt{repOk()} con las condiciones de validez pedidas.
  \item Se corrigió \texttt{pop()} para que:
  \begin{itemize}
    \item devuelva el elemento removido del tope,
    \item limpie la celda desapilada (\texttt{elems[sp] = null}).
  \end{itemize}
  \item Se incorporó validación en \texttt{push} para rechazar \texttt{null}, en consistencia con el invariante.
\end{enumerate}

\subsection{Inconvenientes encontrados}
\begin{itemize}
  \item Para validar estados internos inválidos de \texttt{repOk} (por ejemplo, \texttt{sp} fuera de rango o huecos internos), la API pública no fue suficiente y se utilizó reflexión en los tests.
  \item Al implementar \texttt{repOk} se detectó una inconsistencia adicional en \texttt{pop}: no limpiaba la posición desapilada y devolvía el nuevo tope en vez del elemento removido.
\end{itemize}

\subsection{Ejecución}
\begin{lstlisting}[style=codestyle]
mvn -Dmaven.repo.local=.m2 -Djacoco.skip=true -Dtest=StackArTest test
\end{lstlisting}

\section{Ejercicio 4}
Se implementaron tests parametrizados en CSV para el método:
\begin{itemize}
  \item \texttt{Min.min(List<? extends T>)}
\end{itemize}

\subsection{Archivos}
\begin{itemize}
  \item \texttt{src/test/java/assignment2\_exercises/MinCsvParameterizedTest.java}
  \item \texttt{src/test/resources/assignment2\_exercises/min\_valid\_cases.csv}
  \item \texttt{src/test/resources/assignment2\_exercises/min\_invalid\_cases.csv}
\end{itemize}

\subsection{Cobertura de casos}
\begin{itemize}
  \item Casos válidos con enteros: mínimo al inicio, al medio, al final, repetidos, negativos y listas unitarias.
  \item Casos inválidos desde CSV:
  \begin{itemize}
    \item lista vacía (\texttt{IllegalArgumentException}),
    \item lista con \texttt{null} en primera posición (\texttt{NullPointerException}),
    \item lista con \texttt{null} en otra posición (\texttt{NullPointerException}).
  \end{itemize}
  \item Casos negativos adicionales:
  \begin{itemize}
    \item lista \texttt{null} (\texttt{NullPointerException}),
    \item elementos no mutuamente comparables (\texttt{ClassCastException}).
  \end{itemize}
\end{itemize}

\subsection{Ejecución}
\begin{lstlisting}[style=codestyle]
mvn -Dmaven.repo.local=.m2 -Djacoco.skip=true -Dtest=MinCsvParameterizedTest test
\end{lstlisting}

\section{Ejercicio 5}
Se construyeron tests parametrizados para \texttt{ZuneBug.currentYear(int days)} y, a partir de los fallos observados, se realizó debugging y corrección del método.

\subsection{Archivos}
\begin{itemize}
  \item \texttt{src/test/java/assignment2\_exercises/ZuneBugParameterizedTest.java}
  \item \texttt{src/main/java/assignment2\_exercises/ZuneBug.java}
  \item \texttt{ejercicio5/currentYear\_original.java.txt} (versión previa a la corrección)
\end{itemize}

\subsection{Cobertura de pruebas}
Se agregaron dos grupos parametrizados:
\begin{itemize}
  \item casos de borde (por ejemplo, \texttt{365}, \texttt{366}, \texttt{367}, \texttt{731}, \texttt{1461}),
  \item múltiples valores para comparar \texttt{currentYear} con una función \textit{oracle}.
\end{itemize}

También se utilizó \texttt{assertTimeoutPreemptively} para detectar cuelgues.

\subsection{Fallas detectadas antes de corregir}
\begin{itemize}
  \item Con \texttt{days = 366}, el método podía entrar en bucle infinito.
  \item En cambios exactos de año (por ejemplo, \texttt{731} y \texttt{1461}), devolvía un año menos.
\end{itemize}

\subsection{Causa del defecto}
La lógica original mezclaba condiciones \texttt{> 365} y \texttt{> 366}, con tratamiento incorrecto de límites exactos.

\subsection{Corrección aplicada}
Se reescribió el bucle utilizando \texttt{daysInYear} del año actual:
\begin{itemize}
  \item si \texttt{days >= daysInYear}, se descuenta ese año y se avanza al siguiente,
  \item en caso contrario, se finaliza y se retorna el año actual.
\end{itemize}

Esta corrección elimina tanto el bucle infinito como el error \textit{off-by-one} en los bordes.

\subsection{Ejecución y resultado}
\begin{lstlisting}[style=codestyle]
mvn -Dmaven.repo.local=.m2 -Djacoco.skip=true -Dtest=ZuneBugParameterizedTest test
\end{lstlisting}

Resultado reportado luego de corregir:
\begin{itemize}
  \item \texttt{Tests run: 21, Failures: 0, Errors: 0}
  \item \texttt{BUILD SUCCESS}
\end{itemize}

\section{Ejercicio 6}
Se implementó \texttt{repOK()} en \texttt{BoundedQueue} y se diseñaron tests parametrizados para verificar que el invariante de representación se mantenga luego de operaciones fundamentales de cola.

\subsection{Archivos}
\begin{itemize}
  \item \texttt{src/main/java/assignment2\_exercises/queue/BoundedQueue.java}
  \item \texttt{src/test/java/assignment2\_exercises/queue/BoundedQueueParameterizedTest.java}
\end{itemize}

\subsection{Invariante validado por \texttt{repOK}}
La implementación verifica, entre otros aspectos:
\begin{itemize}
  \item consistencia estructural básica (\texttt{elements}, \texttt{capacity}, \texttt{size}),
  \item rangos válidos de \texttt{front} y \texttt{back},
  \item relación de cola circular: \texttt{back == (front + size) \% capacity},
  \item ausencia de \texttt{null} en la zona ocupada,
  \item ausencia de basura fuera de la zona ocupada (celdas en \texttt{null}).
\end{itemize}

\subsection{Tests parametrizados}
Se utilizó \texttt{@MethodSource} con dos grupos:
\begin{itemize}
  \item escenarios válidos: secuencias de \texttt{enQueue/deQueue} (incluyendo \textit{wrap-around}) y chequeo de \texttt{repOK} tras cada operación,
  \item escenarios inválidos: operaciones que deben lanzar excepción (\texttt{deQueue} en vacía, \texttt{enQueue(null)}, \texttt{enQueue} en llena) y verificación de que \texttt{repOK} permanezca en \texttt{true} después del error.
\end{itemize}

\subsection{Ejecución}
\begin{lstlisting}[style=codestyle]
mvn -Dmaven.repo.local=.m2 -Djacoco.skip=true -Dtest=BoundedQueueParameterizedTest test
\end{lstlisting}

\section*{Conclusiones}
\addcontentsline{toc}{section}{Conclusiones}
La resolución del Práctico 2 consolidó el uso de pruebas parametrizadas, análisis de invariantes de representación y depuración de defectos reales. Los ejercicios permitieron integrar técnicas de diseño de tests con verificación de comportamiento funcional, validación estructural y corrección de implementaciones en Java.

\end{document}
